
\documentclass{article}
\usepackage{colortbl}
\usepackage{makecell}
\usepackage{multirow}
\usepackage{supertabular}

\begin{document}

\newcounter{utterance}

\centering \large Interaction Transcript for game `eqbench', experiment `eq\_bench', episode 59 with reflexion{-}llama{-}3.1{-}8b{-}instruct.
\vspace{24pt}

{ \footnotesize  \setcounter{utterance}{1}
\setlength{\tabcolsep}{0pt}
\begin{supertabular}{c@{$\;$}|p{.15\linewidth}@{}p{.15\linewidth}p{.15\linewidth}p{.15\linewidth}p{.15\linewidth}p{.15\linewidth}}
    \# & \multicolumn{2}{c}{Player} && \multicolumn{2}{c}{Game Master} \\
    \hline

    \theutterance \stepcounter{utterance}  
    & & & \multicolumn{4}{p{0.6\linewidth}}{
        \cellcolor[rgb]{0.9,0.9,0.9}{
            \makecell[{{p{\linewidth}}}]{
                \texttt{\tiny{[P1$\langle$GM]}}
                \texttt{You have attempted this same task before and did not succeed. These are the plans you wrote after those attempts.} \\
\texttt{Trial \#1: New plan: } \\
\\ 
\texttt{Upon re{-}reading the dialogue, I realized that I focused too much on the surface{-}level emotions expressed by the characters, but didn't consider the underlying dynamics of their relationship. Claudia's attempts to reconnect with Tyler are met with resistance and accusations of abandonment, which suggests that Claudia's feelings are more complex than just embarrassment or remorse.} \\
\\ 
\texttt{In my next attempt, I will consider the following:} \\
\\ 
\texttt{1. Claudia's actions and words are perceived as hurtful and manipulative by Tyler, which suggests that Claudia may be feeling a sense of desperation to repair their relationship, but also a sense of guilt and responsibility for the hurt she caused.} \\
\texttt{2. Claudia's repeated attempts to justify her actions and deny Tyler's accusations suggest that she may be feeling a sense of defensiveness and embarrassment, but also a sense of remorse for the pain she caused.} \\
\texttt{3. The fact that Claudia is willing to confront Tyler and face his anger suggests that she may be feeling a sense of determination and desperation to make things right, but also a sense of vulnerability and openness to Tyler's perspective.} \\
\\ 
\texttt{With these considerations in mind, I will revise my scores as follows:} \\
\\ 
\texttt{Embarrassment: 6} \\
\texttt{Desperation: 8} \\
\texttt{Indifference: 0} \\
\texttt{Remorse: 9} \\
\\ 
\texttt{Trial \#2: New plan: } \\
\\ 
\texttt{Upon re{-}reading the dialogue, I realized that I focused too much on the surface{-}level emotions expressed by the characters, but didn't consider the underlying dynamics of their relationship. Claudia's attempts to reconnect with Tyler are met with resistance and accusations of abandonment, which suggests that Claudia's feelings are more complex than just embarrassment or remorse.} \\
\\ 
\texttt{In my next attempt, I will consider the following:} \\
\\ 
\texttt{1. Claudia's actions and words are perceived as hurtful and manipulative by Tyler, which suggests that Claudia may be feeling a sense of desperation to repair their relationship, but also a sense of guilt and responsibility for the hurt she caused.} \\
\texttt{2. Claudia's repeated attempts to justify her actions and deny Tyler's accusations suggest that she may be feeling a sense of defensiveness and embarrassment, but also a sense of remorse for the pain she caused.} \\
\texttt{3. The fact that Claudia is willing to confront Tyler and face his anger suggests that she may be feeling a sense of determination and desperation to make things right, but also a sense of vulnerability and openness to Tyler's perspective.} \\
\\ 
\texttt{With these considerations in mind, I will revise my scores as follows:} \\
\\ 
\texttt{Embarrassment: 6} \\
\texttt{Desperation: 8} \\
\texttt{Indifference: 0} \\
\texttt{Remorse: 9} \\
\\ 
\\ 
\\ 
\texttt{Your task is to predict the likely emotional responses of a character in this dialogue:} \\
\\ 
\texttt{Claudia: Tyler, you look well. Teaching suits you.} \\
\texttt{Tyler: Yeah, well, someone's gotta do the dirty work.} \\
\texttt{Claudia: I didn't mean it like that. I...} \\
\texttt{Tyler: I know what you meant. You always do, don't you?} \\
\texttt{Claudia: Not always, no. But I try to understand, at least.} \\
\texttt{Tyler: Understand? Like you understood when you left me alone?} \\
\texttt{Claudia: That's not fair, Ty. I had to...} \\
\texttt{Tyler: What? Chase your dreams? Leave your past behind?} \\
\texttt{Claudia: I never left you behind, Tyler. You were always...} \\
\texttt{Tyler: Always what? An afterthought? A burden?} \\
\texttt{Claudia: No, never a burden. You were my responsibility.} \\
\texttt{Tyler: Responsibility? Is that why you're here now? After all these years?} \\
\texttt{Claudia: I'm here because I...I wanted to see you.} \\
\texttt{Tyler: Or maybe you're here because you need something. Like always.} \\
\texttt{Claudia: That's not true, Tyler. I...} \\
\texttt{Tyler: Save it, Claudia. I've heard enough of your lies.} \\
\texttt{Claudia: I'm not lying, Tyler. I...} \\
\texttt{Tyler: You abandoned me, Claudia. You can't just waltz back into my life and pretend everything is okay.} \\
\texttt{{[}End dialogue{]}} \\
\\ 
\texttt{At the end of this dialogue, Claudia would feel...} \\
\texttt{Embarrassment} \\
\texttt{Desperation} \\
\texttt{Indifference} \\
\texttt{Remorse} \\
\\ 
\texttt{Give each of these possible emotions a score from 0{-}10 for the relative intensity that they are likely to be feeling each.} \\
\\ 
\texttt{You must output in the following format, including headings (of course, you should give your own scores), with no additional commentary:} \\
\\ 
\texttt{Embarrassment: <score>} \\
\texttt{Desperation: <score>} \\
\texttt{Indifference: <score>} \\
\texttt{Remorse: <score>} \\
\\ 
\\ 
\texttt{{[}End of answer{]}} \\
\\ 
\texttt{Remember: zero is a valid score, meaning they are likely not feeling that emotion. You must score at least one emotion > 0.} \\
\\ 
\texttt{Your answer:} \\
            }
        }
    }
     \\ \\

    \theutterance \stepcounter{utterance}  
    & \multicolumn{4}{p{0.6\linewidth}}{
        \cellcolor[rgb]{0.561,0.737,0.561}{
            \makecell[{{p{\linewidth}}}]{
                \texttt{\tiny{[P1$\rangle$GM]}}
                \texttt{Embarrassment: 8} \\
\texttt{Desperation: 9} \\
\texttt{Indifference: 0} \\
\texttt{Remorse: 7} \\
            }
        }
    }
    & & \\ \\

    \theutterance \stepcounter{utterance}  
    & & & \multicolumn{2}{p{0.3\linewidth}}{
        \cellcolor[rgb]{0.9,0.9,0.9}{
            \makecell[{{p{\linewidth}}}]{
                \texttt{\tiny{[GM$|$GM]}}
                \texttt{Embarrassment: 8} \\
\texttt{Desperation: 9} \\
\texttt{Indifference: 0} \\
\texttt{Remorse: 7} \\
            }
        }
    }
    & & \\ \\

    \theutterance \stepcounter{utterance}  
    & & & \multicolumn{2}{p{0.3\linewidth}}{
        \cellcolor[rgb]{0.9,0.9,0.9}{
            \makecell[{{p{\linewidth}}}]{
                \texttt{\tiny{[GM$|$GM]}}
                \texttt{Embarrassment: 6} \\
\texttt{Desperation: 7} \\
\texttt{Indifference: 0} \\
\texttt{Remorse: 8} \\
            }
        }
    }
    & & \\ \\

    \theutterance \stepcounter{utterance}  
    & & & \multicolumn{2}{p{0.3\linewidth}}{
        \cellcolor[rgb]{0.9,0.9,0.9}{
            \makecell[{{p{\linewidth}}}]{
                \texttt{\tiny{[GM$|$GM]}}
                \texttt{game\_result = LOSE} \\
            }
        }
    }
    & & \\ \\

    \theutterance \stepcounter{utterance}  
    & & \multicolumn{4}{p{0.6\linewidth}}{
        \cellcolor[rgb]{0.95,0.95,0.95}{
            \makecell[{{p{\linewidth}}}]{
                \texttt{\tiny{[P1$\leftrightarrow$P]}}
                \texttt{Post{-}episode Reflexion (not sent to GM):} \\
\texttt{New plan: } \\
\\ 
\texttt{Upon re{-}reading the dialogue, I realized that I focused too much on the surface{-}level emotions expressed by the characters, but didn't consider the underlying dynamics of their relationship. Claudia's attempts to reconnect with Tyler are met with resistance and accusations of abandonment, which suggests that Claudia's feelings are more complex than just embarrassment or remorse.} \\
\\ 
\texttt{In my next attempt, I will consider the following:} \\
\\ 
\texttt{1. Claudia's actions and words are perceived as hurtful and manipulative by Tyler, which suggests that Claudia may be feeling a sense of desperation to repair their relationship, but also a sense of guilt and responsibility for the hurt she caused.} \\
\texttt{2. Claudia's repeated attempts to justify her actions and deny Tyler's accusations suggest that she may be feeling a sense of defensiveness and embarrassment, but also a sense of remorse for the pain she caused.} \\
\texttt{3. The fact that Claudia is willing to confront Tyler and face his anger suggests that she may be feeling a sense of determination and desperation to make things right, but also a sense of vulnerability and openness to Tyler's perspective.} \\
\\ 
\texttt{However, I also realized that I didn't consider the fact that Claudia's attempts to reconnect with Tyler are met with resistance and accusations of abandonment, which suggests that Claudia's feelings may be more complex than just desperation or remorse. I will revise my scores as follows:} \\
\\ 
\texttt{Embarrassment: 7} \\
\texttt{Desperation: 9} \\
\texttt{Indifference: 0} \\
\texttt{Remorse: 8} \\
\\ 
\texttt{In my next attempt, I will also consider the following:} \\
\\ 
\texttt{* Claudia's tone and language may be perceived as insincere or manipulative by Tyler, which could increase her sense of desperation and remorse.} \\
\texttt{* The fact that Claudia is willing to confront Tyler and face his anger suggests that she may be feeling a sense of determination and desperation to make things right, but also a sense of vulnerability and openness to Tyler's perspective.} \\
\texttt{* The underlying dynamics of their relationship, including Claudia's actions and words, may be contributing to Tyler's resistance and accusations of abandonment, which could increase Claudia's sense of guilt and responsibility for the hurt she caused.} \\
            }
        }
    }
    & \\ \\

\end{supertabular}
}

\end{document}
